% Part: first-order-logic
% Chapter: completeness
% Section: construction-of-model

\documentclass[../../../include/open-logic-section]{subfiles}

\begin{document}

\iftag{FOL}
      {\olfileid{fol}{com}{mod}}
      {\olfileid{pl}{com}{mod}}

\olsection{模型的构造}

\begin{explain}
\iftag{FOL}{目前我们暂不关心 $\eq$，即我们只希望证明：一个不含~$\eq$ 的一致语句集~$\Gamma$ 是可满足的。我们首先将~$\Gamma$ 扩充为一个一致的、完备的、饱和的语句集~$\Gamma^*$。在这种情况下，模型~$\Struct{M(\Gamma^*)}$ 的定义十分简单：我们取语言~$\Lang{L'}$ 的闭项的集合作为论域。我们将每个常项符号指派给它自身，并更一般地确保，对每个闭项~$t$，$\Value{t}{M(\Gamma^*)} = t$。谓词符号被指派外延，使得一个原子语句在 $\Struct{M(\Gamma^*)}$ 中为真当且仅当它属于~$\Gamma^*$。这显然使得 $\Gamma^*$ 中所有的原子语句在 $\Struct{M(\Gamma^*)}$ 中为真。只要我们起始的 $\Gamma^*$ 是一致的、完备的且饱和的，其余语句亦将为真。}{我们现在准备定义一个使所有 $!A \in \Gamma$ 为真的赋值。为此，我们首先应用 Lindenbaum（林德鲍姆）引理：我们得到一个完备且一致的 $\Gamma^* \supseteq \Gamma$。我们让 $\Gamma^*$ 中的命题变元来确定 $\pAssign v(\Gamma^*)$。}
\end{explain}

\iftag{FOL}{%


\begin{defn}[项模型]
\ollabel{defn:termmodel}
令 $\Gamma^*$ 为语言~$\Lang L$ 中一个完备、一致且饱和的语句集。$\Gamma^*$ 的\emph{项模型}~$\Struct M(\Gamma^*)$ 定义如下：
\begin{enumerate}
\item 论域~$\Domain{M(\Gamma^*)}$ 是 $\Lang L$ 的所有闭项的集合。
\item 常项符号 $c$ 的解释就是 $c$ 本身：$\Assign{c}{M(\Gamma^*)} = c$。
\item 函数符号 $f$ 被指派为这样一个函数：给定闭项 $t_1$, \dots, $t_n$ 作为自变元，其值为闭项 $f(t_1, \dots, t_n)$：
\[
\Assign{f}{M(\Gamma^*)}(t_1, \dots, t_n) = f(t_1,\dots, t_n)
\]
\item 若 $R$ 是一个 $n$ 元谓词符号，则
  \[
  \tuple{t_1, \dots,
  t_n} \in \Assign{R}{M(\Gamma^*)} \text{ 当且仅当 } \Atom{R}{t_1, \dots,
    t_n} \in \Gamma^*.
  \]
\end{enumerate}
\end{defn}
}{%
\begin{defn}
\ollabel{defn:termmodel}
设 $\Gamma^*$ 是一个完备且一致的公式集。那么我们令
\[
\pAssign v(\Gamma^*)(p) = \begin{cases}
  \True & \text{若 $p \in \Gamma^*$}\\
  \False & \text{若 $p \notin \Gamma^*$}
  \end{cases}
\]
\end{defn}
}

\iftag{FOL}{%
我们现在来验证我们确实有 $\Value{t}{M(\Gamma^*)} = t$。

\begin{lem}
\ollabel{lem:val-in-termmodel} 令 $\Struct M(\Gamma^*)$ 为 \olref{defn:termmodel} 的项模型，则 $\Value{t}{M(\Gamma^*)} = t$。
\end{lem}
\begin{proof}
 对 $t$ 施归纳来证明。基始步（$t$ 是常项符号的情形）直接由项模型的定义得出。对于归纳步，假设 $t_1, \ldots, t_n$ 是闭项，使得 $\Value{t_i}{M(\Gamma^*)} = t_i$，并且 $f$ 是一个 $n$ 元函数符号。那么
\begin{align*}
\Value{f(t_1,\ldots,t_n)}{M(\Gamma^*)} &= \Assign{f}{M(\Gamma^*)}(\Value{t_1}
 {M(\Gamma^*)},
\ldots, \Value{t_n}{M(\Gamma^*)}) \\
&= \Assign{f}{M(\Gamma^*)}(t_1, \dots, t_n) \\
&= f(t_1,\dots, t_n),
\end{align*}
从而由归纳可知，这对每个闭项~$t$ 都成立。
\end{proof}
}{}

\iftag{FOL}{%
\begin{explain}
  \iftag{prvEx}{一个结构 $\Struct{M}$ 可能使一个存在量化语句 $\lexists[x][!A(x)]$ 为真，却没有任何实例 $!A(t)$ 在其中为真。}{}
  \iftag{prvAll}{一个结构 $\Struct{M}$ 可能使一个全称量化语句 $\lforall[x][!A(x)]$ 的所有实例 $!A(t)$ 都为真，却不使 $\lforall[x][!A(x)]$ 本身为真。}{} 这是因为一般而言，并非 $\Domain{M}$ 中的每个元素都是某个闭项的值（$\Struct{M}$ 可能不是覆盖结构）。这就是为什么满足关系要通过变元指派来定义。然而，对于我们的项模型 $\Struct{M(\Gamma^*)}$，这就不必要了——因为它是覆盖结构。这正是下一个结果的内容。
\end{explain}

\begin{prop}
\ollabel{prop:quant-termmodel}
令 $\Struct M(\Gamma^*)$ 为 \olref{defn:termmodel} 的项模型。
\begin{tagenumerate}{prvEx,prvAll}
\tagitem{prvEx}{$\Sat{M(\Gamma^*)}{\lexists[x][!A(x)]}$ 当且仅当存在至少一个闭项 $t$ 使得 $\Sat{M(\Gamma^*)}{!A(t)}$。}{}
\tagitem{prvAll}{$\Sat{M(\Gamma^*)}{\lforall[x][!A(x)]}$ 当且仅当对所有闭项 $t$，$\Sat{M(\Gamma^*)}{!A(t)}$。}{}
\end{tagenumerate}
\end{prop}

\begin{proof}
\begin{tagenumerate}{prvEx,prvAll}
\tagitem{prvEx}{%
  \iftag{probEx}{练习。}{根据\olref[syn][ass]{prop:sat-quant}，
    $\Sat{M(\Gamma^*)}{\lexists[x][!A(x)]}$ 当且仅当存在至少一个变元指派~$s$ 使得 $\Sat{M(\Gamma^*)}{!A(x)}[s]$。
    由于 $\Domain{M(\Gamma^*)}$ 由~$\Lang{L}$ 的闭项组成，
    这当且仅当存在至少一个闭项~$t$ 使得 $s(x) = t$ 且 $\Sat{M(\Gamma^*)}{!A(x)}[s]$。
    由\olref[fol][syn][ext]{prop:ext-formulas}，
    在 $s(x) = t$ 时，$\Sat{M(\Gamma^*)}{!A(x)}[s]$ 当且仅当 $\Sat{M(\Gamma^*)}{!A(t)}[s]$。
    又根据\olref[fol][syn][ass]{prop:sentence-sat-true}，
    由于 $!A(t)$ 是语句，$\Sat{M(\Gamma^*)}{!A(t)}[s]$ 当且仅当 $\Sat{M(\Gamma^*)}{!A(t)}$。}}{}
\tagitem{prvAll}{%
  \iftag{probAll}{练习。}{根据\olref[syn][ass]{prop:sat-quant}，
    $\Sat{M(\Gamma^*)}{\lforall[x][!A(x)]}$ 当且仅当对任意变元指派 $s$，$\Sat{M(\Gamma^*)}{!A(x)}[s]$。
    回顾 $\Domain{M(\Gamma^*)}$ 由~$\Lang{L}$ 的闭项组成，
    因此对每个闭项~$t$，$s(x) = t$ 就是这样的变元指派，
    且对任意变元指派，$s(x)$ 都是某个闭项~$t$。
    由\olref[fol][syn][ext]{prop:ext-formulas}，
    在 $s(x) = t$ 时，$\Sat{M(\Gamma^*)}{!A(x)}[s]$ 当且仅当 $\Sat{M(\Gamma^*)}{!A(t)}[s]$。
    又根据\olref[fol][syn][ass]{prop:sentence-sat-true}，
    由于 $!A(t)$ 是语句，$\Sat{M(\Gamma^*)}{!A(t)}[s]$ 当且仅当 $\Sat{M(\Gamma^*)}{!A(t)}$。}}{}
\end{tagenumerate}
\end{proof}
}{}

\tagprob[FOL]{probEx,probAll}
\begin{prob}
完成命题 \olref[fol][com][mod]{prop:quant-termmodel} 的证明。
\end{prob}
\tagendprob

\begin{lem}[真引理]
\ollabel{lem:truth} \iftag{FOL}{设 $!A$ 不包含~$\eq$。则}{}
$\iftag{FOL}{\Sat{M(\Gamma^*)}}{\pSat{v(\Gamma^*)}}{!A}$ 当且仅当 $!A \in \Gamma^*$。
\end{lem}

\begin{proof}
我们同时证明两个方向，并对 $!A$ 施归纳。
\begin{enumerate}
\tagitem{prvFalse}{\indcase{!A}{\lfalse}
  {根据满足的定义，$\iftag{FOL}{\Sat/{M(\Gamma^*)}{\lfalse}}{\pSat/{v(\Gamma^*)}{\lfalse}}$。
    另一方面，由于 $\Gamma^*$ 一致，故 $\lfalse \notin
    \Gamma^*$。}}{}

\tagitem{prvTrue}{\indcase{!A}{\ltrue}
  {根据满足的定义，$\iftag{FOL}{\Sat{M(\Gamma^*)}}{\pSat{v(\Gamma^*)}}{\ltrue}$。
    另一方面，由于 $\Gamma^*$ 完备且一致，且
    $\Gamma^* \Proves \ltrue$，故 $\ltrue \in \Gamma^*$。}}{}

\tagitem{FOL}{\indcase{!A}{R(t_1, \dots, t_n)}
  {$\Sat{M(\Gamma^*)}{\Atom{R}{t_1, \dots, t_n}}$ 当且仅当 $\tuple{t_1,
      \dots, t_n} \in \Assign{R}{M(\Gamma^*)}$（根据满足的定义），
    当且仅当 $R(t_1, \dots, t_n) \in \Gamma^*$（根据 $\Struct
    M(\Gamma^*)$ 的构造）。}}{\indcase{!A}{p}{$\pSat{v(\Gamma^*)}{p}$ 当且仅当
    $\pAssign v(\Gamma^*)(p) = \True$（根据满足的定义），
    当且仅当 $p \in \Gamma^*$（根据 $\pAssign v(\Gamma^*)$ 的构造）。}}

\tagitem{prvNot}{%
  \indcase{!A}{\lnot !B}
    {$\iftag{FOL}{\Sat{M(\Gamma^*)}}{\pSat{v(\Gamma^*)}}{\indfrm}$ 当且仅当
    $\iftag{FOL}{\Sat/{M(\Gamma^*)}{!B}}{\pSat/{v(\Gamma^*)}{!B}}$（根据满足的定义）。
    由归纳假设，
    $\iftag{FOL}{\Sat/{M(\Gamma^*)}{!B}}{\pSat/{v(\Gamma^*)}{!B}}$ 当且仅当 $!B \notin \Gamma^*$。
    由于 $\Gamma^*$ 一致且完备，$!B \notin \Gamma^*$ 当且仅当 $\lnot !B \in \Gamma^*$。}}{}

\tagitem{prvAnd}{%
\iftag{probAnd}{%
  \indcase!{!A}{!B \land !C}{}}{%
  \indcase{!A}{!B \land
    !C}{$\iftag{FOL}{\Sat{M(\Gamma^*)}}{\pSat{v(\Gamma^*)}}{\indfrm}$
    当且仅当 $\iftag{FOL}{\Sat{M(\Gamma^*)}}{\pSat{v(\Gamma^*)}}{!B}$ 且
    $\iftag{FOL}{\Sat{M(\Gamma^*)}}{\pSat{v(\Gamma^*)}}{!C}$ 均成立（根据满足的定义），当且仅当 $!B \in \Gamma^*$ 且 $!C \in
    \Gamma^*$（由归纳假设）。根据
    \olref[ccs]{prop:ccs}\olref[ccs]{prop:ccs-and}，这当且仅当 $(!B \land !C) \in \Gamma^*$。}}}{}

\tagitem{prvOr}{%
\iftag{probOr}{%
  \indcase!{!A}{!B \lor !C}{}}{%
  \indcase{!A}{!B \lor !C}
    {$\iftag{FOL}{\Sat{M(\Gamma^*)}}{\pSat{v(\Gamma^*)}}{\indfrm}$
    当且仅当 $\iftag{FOL}{\Sat{M(\Gamma^*)}}{\pSat{v(\Gamma^*)}}{!B}$ 或
    $\iftag{FOL}{\Sat{M(\Gamma^*)}}{\pSat{v(\Gamma^*)}}{!C}$（根据满足的定义），当且仅当 $!B \in \Gamma^*$ 或 $!C \in
    \Gamma^*$（由归纳假设）。根据
    \olref[ccs]{prop:ccs}\olref[ccs]{prop:ccs-or}，这当且仅当 $(!B
    \lor !C) \in \Gamma^*$。}}}{}

\tagitem{prvIf}{%
\iftag{probIf}{%
  \indcase!{!A}{!B \lif !C}{}}{%
  \indcase{!A}{!B \lif !C}
    {$\iftag{FOL}{\Sat{M(\Gamma^*)}}{\pSat{v(\Gamma^*)}}{\indfrm}$
    当且仅当 $\iftag{FOL}{\Sat/{M(\Gamma^*)}}{\pSat/{v(\Gamma^*)}}{!B}$ 或 $\iftag{FOL}{\Sat{M (\Gamma^*)}}{\pSat{v(\Gamma^*)}}{!C}$（根据满足的定义），当且仅当 $!B \notin \Gamma^*$ 或 $!C \in
    \Gamma^*$（由归纳假设）。根据
    \olref[ccs]{prop:ccs}\olref[ccs]{prop:ccs-if}，这当且仅当 $(!B
    \lif !C) \in \Gamma^*$。}}}{}

\iftag{FOL}{%
\tagitem{prvAll}{%
  \iftag{probAll}{%
    \indcase!{!A}{\lforall[x][!B(x)]}{}}{%
    \indcase{!A}{\lforall[x][!B(x)]}{$\Sat{M(\Gamma^*)}{\indfrm}$ 当且仅当
      对所有项~$t$，$\Sat{M(\Gamma^*)}{!B(t)}$ 成立（\olref{prop:quant-termmodel}）。
      由归纳假设，这当且仅当对所有项~$t$，$!B(t) \in \Gamma^*$；而根据
      \olref[hen]{prop:saturated-instances}，这又当且仅当 $\lforall[x][!A(x)] \in \Gamma^*$。}}}{}

\tagitem{prvEx}{%
  \iftag{probEx}{%
    \indcase!{!A}{\lexists[x][!B(x)]}{}}{%
    \indcase{!A}{\lexists[x][!B(x)]}{$\Sat{M(\Gamma^*)}{\indfrm}$ 当且仅当
      存在至少一个项~$t$ 使得 $\Sat{M(\Gamma^*)}{!B(t)}$ 成立（\olref{prop:quant-termmodel}）。
      由归纳假设，这当且仅当存在至少一个项~$t$ 使得 $!B(t) \in \Gamma^*$。
      根据 \olref[hen]{prop:saturated-instances}，这又当且仅当 $\lexists[x][!B(x)] \in \Gamma^*$。}}}{}
}{}
\end{enumerate}
\end{proof}


\tagprob[FOL]{probOr,probAnd,probIf,probEx,probAll}

\begin{prob}
完成 \olref[fol][com][mod]{lem:truth} 的证明。
\end{prob}

\tagendprob

\tagprob[notFOL]{probOr,probAnd,probIf}

\begin{prob}
完成 \olref[pl][com][mod]{lem:truth} 的证明。
\end{prob}

\tagendprob

\end{document}